\section{Exact contraction on the open square grid}

\begin{definition}[Forward square edges]
    \label{def:peps_forward_square_edges}
    \lean{TNLean.PEPS.Approximation.Vertex,TNLean.PEPS.Approximation.ForwardAdjacent,TNLean.PEPS.Approximation.ForwardEdge}
    \leanok
    For $\Lambda=\{0,\ldots,L-1\}^2$, a forward edge is an ordered pair
    $((x,y),(x+1,y))$ or $((x,y),(x,y+1))$ with both endpoints in $\Lambda$.
\end{definition}

\begin{definition}[Square-edge correspondence]
    \label{def:peps_square_edge_correspondence}
    \lean{TNLean.PEPS.Approximation.forwardSquareEdgeEquiv,TNLean.PEPS.Approximation.forwardSquareIncidentEquiv,TNLean.PEPS.Approximation.pinnedSquareIncidentEquiv,TNLean.PEPS.Approximation.forwardSquareVirtualConfigEquiv,TNLean.PEPS.Approximation.forwardSquareEdgeEquiv_val,TNLean.PEPS.Approximation.forwardSquareEdgeEquiv_symm_val,TNLean.PEPS.Approximation.forwardSquareIncidentEquiv_val,TNLean.PEPS.Approximation.pinnedSquareIncidentEquiv_val,TNLean.PEPS.Approximation.graphBondDim,TNLean.PEPS.Approximation.graphBondDim_forward,TNLean.PEPS.Approximation.forwardSquareVirtualConfigEquiv_incident,TNLean.PEPS.Approximation.forwardSquareVirtualConfigEquiv_pinnedIncident}
    \leanok
    \uses{def:peps_forward_square_edges}
    The forward edges are in bijection with the lexicographically oriented
    edges of the open square graph by preserving both endpoints.
    For every vertex this gives a bijection of incident edges, and for every
    family of dimensions $(D_e)_e$ it gives a bijection of assignments
    $\prod_e\{0,\ldots,D_e-1\}$, with each coordinate unchanged.
\end{definition}

\begin{theorem}[Equality of square-grid contractions]
    \label{thm:peps_square_grid_contraction}
    \lean{TNLean.PEPS.Approximation.stateCoeff_pinnedTensorToGraphTensor,TNLean.PEPS.Approximation.stateCoeff_vectorTensorToGraphTensor,TNLean.PEPS.Approximation.pinnedTensorToGraphTensor,TNLean.PEPS.Approximation.vectorTensorToGraphTensor,TNLean.PEPS.Approximation.pinnedTensorToGraphTensor_component,TNLean.PEPS.Approximation.vectorTensorToGraphTensor_component}
    \leanok
    \uses{def:peps_square_edge_correspondence}
    Let $A_v$ be arbitrary local tensors with physical dimension $q$ and
    edge-dependent virtual dimensions $D_e$. Transport each virtual coordinate
    along the edge correspondence, and put the physical argument last when it
    was written first. The resulting graph tensor $\widetilde A$ has dimensions
    $\widetilde D_{\iota(e)}=D_e$. For every physical configuration $s$,
    \begin{align}
        \sum_b\prod_{v\in\Lambda}\widetilde A_v(b|_v,s_v)
         & =\sum_a\prod_{v\in\Lambda}A_v(s_v,a|_v).
    \end{align}
    The same equality holds when the source tensor already has its physical
    argument last. The physical-first assertion permits zero virtual dimensions.
\end{theorem}
\begin{proof}
    \leanok
    \uses{def:peps_square_edge_correspondence}
    Reindex the finite sum by the bijection of virtual assignments.
    Its restriction to every incident edge preserves that edge's index, so each
    factor in the vertex product is equal.
\end{proof}


\begin{definition}[Two orders for local tensor arguments]
    \label{def:peps_square_local_presentations}
    \lean{TNLean.PEPS.Approximation.Pinned.IncidentEdge,TNLean.PEPS.Approximation.Pinned.State,TNLean.PEPS.Approximation.Pinned.LocalTensor,TNLean.PEPS.Approximation.Pinned.contractPEPS,TNLean.PEPS.Approximation.Vector.IncidentEdge,TNLean.PEPS.Approximation.Vector.State,TNLean.PEPS.Approximation.Vector.Tensor,TNLean.PEPS.Approximation.Vector.Tensor.contract,TNLean.PEPS.Approximation.Vector.Tensor.maxBondDim}
    \leanok
    \uses{def:peps_forward_square_edges}
    A local tensor may be written as $A_v(s_v,a|_v)$ or $A_v(a|_v,s_v)$.
    In either case its contraction is the vector in
    $\ell^2(\Lambda\to\{0,\ldots,q-1\};\C)$ whose $s$ coordinate is
    $\sum_a\prod_v A_v(s_v,a|_v)$, with the arguments reversed in the second
    presentation. The second presentation includes the requirement $D_e>0$.
    Its maximum bond dimension is $\max(\{D_e:e\in E\}\cup\{0\})$.
\end{definition}

For a $2\times2$ grid, the contraction has four independent virtual bonds
and four open physical indices. Each bond retains its own dimension.
\begin{tenkzequation}
    \tenkzkernel
    \begin{tenkz}[rows={wire,wire}, cols=2, bonds=none]
        \tn[at=(2,2), name=squareGridTL, skin=box,
            ports={0:virtual,270:virtual,135:physical}]{A_{01}}
        \tn[at=4 e of squareGridTL, name=squareGridTR, skin=box,
            ports={180:virtual,270:virtual,45:physical}]{A_{11}}
        \tn[at=4 s of squareGridTL, name=squareGridBL, skin=box,
            ports={0:virtual,90:virtual,225:physical}]{A_{00}}
        \tn[at=4 e of squareGridBL, name=squareGridBR, skin=box,
            ports={180:virtual,90:virtual,315:physical}]{A_{10}}
        \tnwire[name=squareGridH1]{squareGridTL.0}{squareGridTR.180}
        \tnmark[form=label, label pos=n]{on squareGridH1 0.5}{$D_{h_1}$}
        \tnwire[name=squareGridH0]{squareGridBL.0}{squareGridBR.180}
        \tnmark[form=label, label pos=s]{on squareGridH0 0.5}{$D_{h_0}$}
        \tnwire[name=squareGridV0]{squareGridTL.270}{squareGridBL.90}
        \tnmark[form=label, label pos=w]{on squareGridV0 0.5}{$D_{v_0}$}
        \tnwire[name=squareGridV1]{squareGridTR.270}{squareGridBR.90}
        \tnmark[form=label, label pos=e]{on squareGridV1 0.5}{$D_{v_1}$}
        \tnwire[name=squareGridPhysical01, via={1 nw of squareGridTL}]{squareGridTL.135}{open nw}
        \tnmark[form=label, label pos=w]
        {on squareGridPhysical01 0.6}{$i_{01}$}
        \tnwire[name=squareGridPhysical11, via={1 ne of squareGridTR}]{squareGridTR.45}{open ne}
        \tnmark[form=label, label pos=e]
        {on squareGridPhysical11 0.6}{$i_{11}$}
        \tnwire[name=squareGridPhysical00, via={1 sw of squareGridBL}]{squareGridBL.225}{open sw}
        \tnmark[form=label, label pos=w]
        {on squareGridPhysical00 0.6}{$i_{00}$}
        \tnwire[name=squareGridPhysical10, via={1 se of squareGridBR}]{squareGridBR.315}{open se}
        \tnmark[form=label, label pos=e]
        {on squareGridPhysical10 0.6}{$i_{10}$}
    \end{tenkz}
\end{tenkzequation}
